\subsection{Related work} \label{sec:relatedwork}
%Paradigms: smoothness, corruptions, ball around a distribution, data-dependent guarantees on regret 
The general problem of combining stochastic and worst-case guarantees has attracted much recent attention---while the optimal algorithms in the stochastic case, such as $\ftl$ and other greedy algorithms, offer the advantages of being easy to implement and working well on typical data~\citep{bastani2021mostly} using these algorithms in high-stakes settings (such as clinical trials) requires stronger guarantees on their performance even under strategic data. 
As discussed previously, there is a large literature on such best-of-both-worlds algorithms for both full and partial information settings~\citep{Wei--Luo2018, Bubeck--Li--Luo--Wei2019, Zimmert--Seldin2019, Dann--Wei--Zimmert2023}, and also more complex settings such as metrical task systems~\citep{bhuyan2023best} and control~\citep{sabag2021regret, goel2023best}. Another line of work~\citep{Rakhlin--Sridharan--Tewari2011, Haghtalab--Roughgarden--Shetty2022, Bhatt--Haghtalab--Shetty2023, Block--Dagan--Golowich--Rakhlin2022} considers \emph{smoothed analysis}, where the worst-case actions of the adversary are perturbed by nature. A third approach~\citep{bubeck2012best,Lykouris--Mirrokni--PaesLeme2018, Amir--Attias--Koren--Mansour--Livni2020, Zimmert--Seldin2019} interpolates between the settings by considering most $\ell_t$ to be i.i.d., interspersed with a few adversarially chosen instances (corruptions). Finally, another line considers the data-generating distribution to come from a ball of specified radius around i.i.d. distributions~\citep{mourtada2019optimality, bilodeau2023relaxing}. While all these approaches provide useful insights into the gap between average and worst-case guarantees, one can argue they are all imprecisely specified -- given an instance $\{\ell_t\}_{t\in[n]}$ in hindsight, there is no clear sense as to which model best `explains' the instance.


Our focus on instance optimality instead follows the approach of better understanding and shaping the per-sequence regret landscape. The origins of this approach arguably come from the seminal work of~\cite{Cover1966} for binary prediction (we discuss this in more detail in~\cref{sec:converse}), with a later focus on better bounds for benign instances in general online learning~\citep{Auer--Cesa-Bianchi--Gentile2002,cesa2005minimizing,hazan2010extracting}. More recently, 
a line of work~\citep{Koolen--VanErven--Grunwald2014LLR, VanErven--Grunwald--Mehta--Reid--Williamson2015FastRates, VanErven--Koolen2016Metagrad, gaillard2014second} have studied designing policies that can adapt to different types of data sequences and achieve multiple performance guarantees simultaneously. The main idea is to use multiple learning rates that are weighted according to their empirical performance on the data. While the focus is still primarily on classifying instances based on when they are easier/harder to learn, some of the resulting guarantees have an instance-optimality flavor; for example,~\cite{VanErven--Koolen2016Metagrad} show how to simultaneously match the performance guarantee (in terms of certain variance bounds) attained by different learning rates in Hedge. Such approaches however need to understand their baseline algorithms in great detail, and use them in a `white-box' way to get their guarantees. 
%~\cite{VanErven--Grunwald--Mehta--Reid--Williamson2015FastRates} bridges the gap between stochastic and worst-case by quantifying ``simpler" sequences to predict via the Bernstein condition, which implies fast learning rates in statistical learning. 

In contrast, our approach fundamentally focuses on combining policies in a black-box way to get instance optimal outcomes. 
As we mention earlier, these all get guarantees with respect to the loss of the reference algorithms, which in general may be much weaker than our regret guarantees, in particular, if done in a way which treats both algorithms symmetrically. One way around this, however, is to tune the learning rates in an \emph{asymmetric} way; this idea is at the heart of the $(\mathcal{A},\mathcal{B})$-PROD algorithm of \cite{sani2014exploiting}, which hedges between the losses of the two algorithms to get a simultaneous  asymmetric
regret guarantee with respect to the two algorithms. In particular, their approach loses at most a constant additional regret compared to FTL; however on sequences where FTL has $\Theta(n)$ regret, $(\mathcal{A},\mathcal{B})$-PROD can only get a regret of $\Theta(\sqrt{n \log n})$, which is not constant factor competitive compared to the $O(\sqrt{n})$ minimax regret.
Similar ideas have been used to balance algorithms in more complex settings with partial information and/or state -- in particular, we note work on corralling bandit algorithms~\cite{agarwal2017corralling,pacchiano2020model,Dann--Wei--Zimmert2023}, as well as online algorithms with predictions~\cite{bamas2020primal,dinitz2022algorithms,anand2022online}.

To our knowledge, the only previous result which attains a comparable instance-optimality guarantee to ours is that of~\cite{de2014follow} for the experts problem, where the authors propose the $\mathsf{FlipFlop}$ policy which interleaves $\mathsf{Hedge}$ (with varying learning rates) and $\mathsf{FTL}$ to obtain a regret guarantee similar to that of Corollary~\ref{cor:SqrtLStarInstantiation}. In fact, their guarantee is stronger as $\mathsf{FlipFlop}$ is shown to be 5.64-competitive with respect to $\min\{\Reg(\ftl,\ell^n), g(L^*)\}$ where $g(L^*) \le \sqrt{L^* \log m}$~\citep[Corollary 16]{de2014follow}. However, 
while $\mathsf{FlipFlop}$ depends on a clever choice of learning rates in $\mathsf{Hedge}$, $\bob$ can black-box interleave $\ftl$ with \emph{any} worst-case/small-loss algorithm without knowing the inner workings of said algorithm, which we see as a significant engineering strength. More importantly, our approach to this problem is distinct as we focus on the fundamental limits (upper and lower bounds) on the \emph{competitive ratio} that must be incurred when combining $\ftl$ with a worst-case policy; to the best of our knowledge this viewpoint, and the corresponding reduction to an optimal stopping problem, has not been previously explored. While our approach is very simple, we believe it is fundamentally different from existing work (in particular, from the approaches behind FlipFlop and $(\mathcal{A},\mathcal{B})$-PROD). In particular, existing approaches are all based on hedging between the \emph{realized losses} of the policies (which are clearly monotone non-decreasing over time); in contrast, our main idea is that for policies like FTL, one can directly construct a monotone estimator for its \emph{regret}, bypassing the need to use loss as a proxy for the regret.
%%%%%%%%%%%%%%%%%%%%%%%%%%%%%%%%%%%%%%%%%%%%%%%%%%%%%%%%%%%%%%%%%%%%%%%%%%%%%%%%%%%%%%%%%%%%%%%%%%%%%%%%%%%%%

